\documentclass[../../main.tex]{subfiles}
\begin{document}


{\bf [解]}

$C$の方程式は
\begin{align}
 C:x^2+y^2=1
\end{align}
である．題意の1より$D$の中心は$C$上にあるから，$P(\cos\theta,\sin\theta,0)$（$0\le\theta<2\pi$）とかける．
題意の2より$D$のある平面は常に$(0,1,0)$すなわち$y$軸と直交するので常に$y=\sin\theta$上にあり，方程式は
\begin{align}
 D: (x-\cos\theta)^2+z^2\le1,\quad y=\sin\theta \label{eq:1}
\end{align}
と書ける．$\theta$を動かした時に\cref{eq:1}の動く領域$S$の体積$V$を求めれば良い．
この図形は$y=0$を軸に対称だから，以下$y \ge 0$の領域のみ考え，その体積を二倍すれば$V$が求まる．


\begin{figure}[htb]
\begin{center}
\begin{tikzpicture}[
    scale=1.8,
    >=stealth,
    % 3D 投影ベクトル (x: 手前斜め左, y: 右, z: 上)
    x={(-0.6cm,-0.35cm)},
    y={(1.0cm,0.0cm)},
    z={(0cm,1.0cm)}
  ]

  % --- 1. パラメータ設定 (theta = 45度の例) ---
  \def\thetaVal{45}
  \pgfmathsetmacro{\cosT}{cos(\thetaVal)}
  \pgfmathsetmacro{\sinT}{sin(\thetaVal)}

  % 中心点 P
  \coordinate (P) at (\cosT, \sinT, 0);

  % --- 2. xy 平面の座標軸 ---
  \draw[->, thick, gray!80] (-1.6,0,0) -- (1.8,0,0) node[below left, black] {$x$};
  \draw[->, thick, gray!80] (0,-1.6,0) -- (0,2.2,0) node[right, black] {$y$};
  \draw[->, thick, gray!80] (0,0,-1.4) -- (0,0,1.8) node[above, black] {$z$};
  \node[below left, font=\scriptsize] at (0,0,0) {$O$};

  % --- 3. xy 平面上の基準円 C: x^2 + y^2 = 1 ---
  \draw[thick, gray!60, domain=0:360, samples=60] 
    plot ({cos(\x)}, {sin(\x)}, 0);
  \node[gray!80!black, font=\scriptsize, below] at (-0.7, -0.7, 0) {$C: x^2+y^2=1$};

  % --- 4. 切断面 y = sin theta (xz 平面に平行な面) ---
  \draw[dashed, blue!40] (-1.8, \sinT, -1.3) -- (1.8, \sinT, -1.3) -- (1.8, \sinT, 1.3) -- (-1.8, \sinT, 1.3) -- cycle;
  \node[blue!70!black, font=\scriptsize, above right] at (1.8, \sinT, 1.3) {平面 $y = \sin\theta$};

  % --- 5. 円板 D: (x - cos theta)^2 + z^2 <= 1 on y = sin theta ---
  \fill[blue!20, opacity=0.7] 
    plot[domain=0:360, samples=60] ({\cosT + cos(\x)}, \sinT, {sin(\x)}) -- cycle;
  \draw[ultra thick, blue!80!black, domain=0:360, samples=60] 
    plot ({\cosT + cos(\x)}, \sinT, {sin(\x)});

  % --- 6. 中心 P とガイド線 ---
  \draw[dashed, red!80!black] (0,0,0) -- (P);
  \fill[red!80!black] (P) circle (1.2pt) node[below right, font=\small] {$P(\cos\theta, \sin\theta, 0)$};

  % 円板 D の半径 1
  \draw[<->, thin, blue!90!black] (P) -- ({\cosT}, \sinT, 1.0)
    node[midway, right, font=\scriptsize] {半径 $1$};

  % 角度 theta
  \draw[->, thin, red!80!black] (0.4, 0, 0) arc (0:\thetaVal:0.4);
  \node[red!80!black, font=\scriptsize] at (0.45, 0.25, 0) {$\theta$};

  % 円板ラベル
  \node[blue!90!black, font=\small] at (\cosT, \sinT, -0.5) {$D$};

\end{tikzpicture}
\end{center}
\caption{立体 $S$ の構成：$xy$ 平面上の円 $C$ に中心をもち $y$ 軸に直交する単位円板 $D$}
\end{figure}

そこで$S$を断面$y=k\, (0 \le k \le 1)$で切った時を考える．$0 \le \theta < 2\pi$より$\sin\theta = k$を満たす$\theta$に対して\cref{eq:1}で対応するのは$\theta,2\pi-\theta$の時で，断面は下図のように二つの円の内側（斜線部）になる．対称性より$0 \le \theta \le \dfrac{\pi}{2}$で考えれば良い．

\begin{figure}[htb]
\begin{center}
\begin{tikzpicture}[scale=2.2, >=stealth]

  % --- 1. パラメータ設定 (theta = 50度) ---
  \def\thetaVal{50}
  \pgfmathsetmacro{\cVal}{cos(\thetaVal)} % cos(theta) ≒ 0.643

  % --- 2. 2円の和集合領域の斜線ハッチング (V(theta)) ---
  \begin{scope}
    \clip (-\cVal, 0) circle (1.0) (\cVal, 0) circle (1.0);
    \fill[pattern=north east lines, pattern color=blue!70] 
      (-1.8, -1.2) rectangle (1.8, 1.2);
  \end{scope}

  % --- 3. 座標軸 (xz 平面) ---
  \draw[->, thick] (-2.0, 0) -- (2.0, 0) node[right] {$x$};
  \draw[->, thick] (0, -1.3) -- (0, 1.4) node[above] {$z$};
  \node[below left, font=\small] at (0,0) {$O$};

  % --- 4. 2つの円の輪郭 ---
  % 右側の円: 中心 (cos theta, 0)
  \draw[ultra thick, blue!80!black] (\cVal, 0) circle (1.0);
  % 左側の円: 中心 (-cos theta, 0)
  \draw[ultra thick, blue!80!black] (-\cVal, 0) circle (1.0);

  % --- 5. 重なり部分（共通部分）の境界強調 ---
  \draw[dashed, red!80!black, thick] 
    plot[domain=-65:65, samples=40] ({-\cVal + cos(\x)}, {sin(\x)});
  \draw[dashed, red!80!black, thick] 
    plot[domain=115:245, samples=40] ({\cVal + cos(\x)}, {sin(\x)});

  % --- 6. 中心・目盛りプロット ---
  \fill[red!80!black] (\cVal, 0) circle (1.2pt) node[below, font=\small] {$\cos\theta$};
  \fill[red!80!black] (-\cVal, 0) circle (1.2pt) node[below, font=\small] {$-\cos\theta$};

  \fill[black] (0, 1) circle (1.0pt) node[above left, font=\scriptsize] {$1$};
  \fill[black] (0, -1) circle (1.0pt) node[below left, font=\scriptsize] {$-1$};

  % 面積ラベル
  \node[blue!90!black, font=\large] at (1.1, 0.4) {$V(\theta)$};
  \node[red!80!black, font=\scriptsize] at (0, 0.4) {共通部分};

\end{tikzpicture}
\end{center}
\caption{断面 $y=k$（$k=\sin\theta$）における断面積 $V(\theta)$}
\end{figure}

この部分の面積$V(\theta)$とすると，これは二つの円の合計の面積$2\pi$から共通部分の面積を引いたもので
\begin{align}
 V(\theta) 
  &= 2\pi - 4\int_{\cos\theta}^{1}\sqrt{1-x^2}\,dx 
\end{align}
である．

\begin{figure}[htb]
\begin{center}
\begin{tikzpicture}[scale=3.5, >=stealth]

  % --- 1. パラメータ設定 (theta = 55度) ---
  \def\thetaVal{55}
  \pgfmathsetmacro{\cVal}{cos(\thetaVal)} % cos(theta) ≒ 0.5735
  \pgfmathsetmacro{\yCut}{sqrt(1 - \cVal*\cVal)} % sin(theta)

  % --- 2. 第 1 象限の積分領域のハッチング (cos theta <= x <= 1, y = sqrt(1-x^2)) ---
  \begin{scope}
    \clip (\cVal, 0) -- plot[domain=\cVal:1.0, samples=40] (\x, {sqrt(1 - \x*\x)}) -- (1.0, 0) -- cycle;
    \fill[pattern=north east lines, pattern color=red!80] 
      (\cVal, 0) rectangle (1.05, 1.05);
  \end{scope}

  % 残りの 3 つの対称領域（薄いハッチング）
  \begin{scope}
    \clip (\cVal, 0) -- plot[domain=\cVal:1.0, samples=40] (\x, {-sqrt(1 - \x*\x)}) -- (1.0, 0) -- cycle;
    \fill[pattern=north west lines, pattern color=red!40] (\cVal, -1.0) rectangle (1.0, 0);
  \end{scope}

  % --- 3. 座標軸 ---
  \draw[->, thick] (-1.2, 0) -- (1.3, 0) node[right] {$x$};
  \draw[->, thick] (0, -1.1) -- (0, 1.2) node[above] {$z$};
  \node[below left, font=\small] at (0,0) {$O$};

  % --- 4. 2つの円弧の描画 ---
  % (x + cos theta)^2 + z^2 = 1 の右半分（原点から見ると x^2+z^2=1 を -cos theta ずらしたもの）
  % 原点を中心として見た単位円弧 x^2 + z^2 = 1
  \draw[ultra thick, blue!80!black, domain=-180:180, samples=50] 
    plot ({cos(\x)}, {sin(\x)}) node[above right, font=\small] {$z = \sqrt{1-x^2}$};

  % --- 5. 積分区間 [cos theta, 1] のガイド線 ---
  \draw[thick, dashed, red!80!black] (\cVal, -\yCut) -- (\cVal, \yCut);
  \fill[red!80!black] (\cVal, 0) circle (0.9pt) node[below left, font=\small] {$\cos\theta$};
  \fill[black] (1.0, 0) circle (0.9pt) node[below right, font=\small] {$1$};

  % 境界点
  \fill[red!80!black] (\cVal, \yCut) circle (1.0pt);
  \fill[red!80!black] (\cVal, -\yCut) circle (1.0pt);

  % --- 6. 積分記号と注記 ---
  \draw[<-, thick, red!80!black] (0.75, 0.3) to[bend left=20] (0.9, 0.75);
  \node[right, font=\small, red!80!black] at (0.9, 0.8) 
    {$\displaystyle \int_{\cos\theta}^1 \sqrt{1-x^2}\,dx$};

  % \node[font=\footnotesize, gray!80!black] at (0, -0.6) {（全体で $4$ 倍）};

\end{tikzpicture}
\end{center}
\caption{共通部分の面積計算}
\end{figure}

$x=\cos\alpha$（$\alpha:\theta\to 0$として$x:c\to 1$）と置換すると，$dx=-\sin\alpha\,d\alpha$，$\sqrt{1-x^2}=\sin\alpha$（$0\le\alpha\le\pi/2$）だから
\begin{align}
 V(\theta)
  &= 2\pi - 4\int_0^\theta\sin^2\alpha\,d\alpha \\
  &= 2\pi - 4\int_0^\theta\left(\dfrac{1-2\cos 2\alpha}{2}\right)\,d\alpha \\
  &= 2\pi - 4\left[\dfrac{\alpha}{2}-\dfrac{\sin2\alpha}{4}\right]_0^{\theta} \\
  &= 2\pi + \sin 2\theta - 2\theta \label{eq:2}
\end{align}
である．

したがって求める体積$V$は\cref{eq:2}を$0\le k=\sin\theta \le 1$で積分した上で$2$倍したもので
\begin{align}
\begin{split}
 \dfrac{V}{2}
   &=\int_0^1V(\theta)\,dk \\
   &=\int_0^{\pi/2}V(\theta)\cdot\dfrac{dk}{d\theta}\,d\theta \\
   &=\int_0^{\pi/2}(2\pi+\sin2\theta-2\theta)\cos\theta\,d\theta \\
   &=\int_0^{\pi/2}\bigl(2\pi\cos\theta+\cos\theta\sin2\theta-2\theta\cos\theta\bigr)d\theta \label{eq:3}
\end{split}
\end{align}
である．
以下各項を計算すると
\begin{align}
\begin{split}
\int_0^{\pi/2}\cos\theta \,d\theta&=\bigl[\sin\theta\bigr]_0^{\pi/2}=1 \\
\int_0^{\pi/2}\theta\cos\theta\,d\theta&=\bigl[\theta\sin\theta+\cos\theta\bigr]_0^{\pi/2}=\left(\dfrac{\pi}{2}+0\right)-(0+1)=\dfrac{\pi}{2}-1 \\
\int_0^{\pi/2}\cos\theta\sin2\theta\,d\theta&=2\int_0^{\pi/2}\sin\theta\cos^2\theta\,d\theta=2\left[-\dfrac{\cos^3\theta}{3}\right]_0^{\pi/2}=\dfrac{2}{3}
\end{split}
\end{align}
だから\cref{eq:3}に代入して
\begin{align}
\begin{split}
\dfrac{V}{2}&=2\pi\cdot1+\dfrac{2}{3}-2\left(\dfrac{\pi}{2}-1\right)=2\pi+\dfrac{2}{3}-\pi+2=\pi+\dfrac{8}{3} \\
\therefore\ 
 V&=2\pi+\dfrac{16}{3}
\end{split}
\end{align}
が求める体積である．


\newpage

\end{document}
